% Generated draft: requires mathematical and source review.
\documentclass[11pt]{article}
\usepackage[margin=25mm]{geometry}
\usepackage{fontspec}
\setmainfont{Latin Modern Roman}
\usepackage{amsmath,amssymb,mathtools,unicode-math}
\setmathfont{Latin Modern Math}
\usepackage{enumitem,needspace,xurl,hyperref}
\setlength{\parindent}{0pt}
\setlength{\parskip}{5pt}
\setlength{\emergencystretch}{3em}
\newcommand{\N}{\mathbb{N}}
\newcommand{\Z}{\mathbb{Z}}
\newcommand{\Q}{\mathbb{Q}}
\newcommand{\R}{\mathbb{R}}
\newcommand{\C}{\mathbb{C}}
\newcommand{\F}{\mathbb{F}}
\newcommand{\A}{\mathbb{A}}
\newcommand{\PP}{\mathbb{P}}
\newcommand{\E}{\mathbb{E}}
\newcommand{\eps}{\varepsilon}
\newcommand{\dd}{\,\mathrm{d}}
\begin{document}
\section*{UnsolvedMath ID: 30003717}
\textbf{Source problem code:} catalogue entry
\subsection{Definitions and mathematical statement}
For a positive integer \(n\), write
\[
\operatorname{Sym}^{+}(n)=\{X\in\R^{n\times n}:X^{\mathsf T}=X,\ v^{\mathsf T}Xv>0\text{ for every }v\in\R^n\setminus\{0\}\}.
\]
Let \(X,Y\in\operatorname{Sym}^{+}(n)\). The source does not specify a dimension or a range of dimensions for the intended assertion. A continuous piecewise-smooth curve joining these matrices can conventionally be represented by \(\gamma:[0,1]\to\operatorname{Sym}^{+}(n)\), with \(\gamma(0)=X\), \(\gamma(1)=Y\), and a finite partition on whose pieces the curve is smooth; derivatives need not agree at partition points. This convention describes the ambient curves but does not supply the missing scaling--rotation structure.

The supplied problem asks whether the scaling--rotation advantage for smooth curves persists for piecewise-smooth curves, and asks for algebraic conditions under which minimal piecewise-smooth scaling--rotation curves are shortest among all such curves.

The following is only a conditional formulation of the shortestness part. Once a comparison class \(\mathcal A(X,Y)\), its length functional \(L\), and the intended class \(\mathcal M_{\mathrm{SR}}(X,Y)\subseteq\mathcal A(X,Y)\) of minimal piecewise-smooth scaling--rotation curves have been specified, characterize algebraically those \(\gamma\in\mathcal M_{\mathrm{SR}}(X,Y)\) satisfying
\[
L(\gamma)=\inf_{\eta\in\mathcal A(X,Y)}L(\eta).
\]
The source does not define these classes, the length functional, the meaning of minimality, or the scaling--rotation advantage. In particular, it does not determine whether the comparison class consists of all piecewise-smooth curves in \(\operatorname{Sym}^{+}(n)\) or only curves with a specified scaling--rotation representation. Consequently, the displayed condition records the intended global shortestness question without supplying missing hypotheses or asserting existence of a minimizer. The advantage question cannot be made precise from the supplied record alone.
\subsection{Short English statement}
For symmetric positive-definite matrices \(X\) and \(Y\), determine whether the stated scaling--rotation advantage extends from smooth to piecewise-smooth curves, and give algebraic criteria for the intended minimal scaling--rotation curves to be globally shortest. The source leaves the advantage, curve classes, length, and meaning of minimality unspecified, so this is a qualified statement-only draft requiring mathematical review.
\subsection{Sources}
\begin{itemize}
\item Local UnsolvedMath registry: supplied complete record for permanent ID \(30003717\), source problem code \texttt{catalogue entry}, title \emph{Shortest Piecewise-Smooth Scaling--Rotation Curves}, in the Oberwolfach Reports Open Problems collection. Repository registry locations are \texttt{data/unsolvedmath/problems.json} and \texttt{catalog/problems.json}; neither was independently inspected for this draft.
\item Supplied report citation: \emph{Statistics for Data with Geometric Structure (2019)}, Oberwolfach Reports, pp.\ 125--186, DOI \texttt{10.4171/owr/2018/3}, \url{https://doi.org/10.4171/owr/2018/3}. The record gives proposed year 2018 and describes the report as documenting workshop open-problem discussions. No individual proposer is supplied.
\item Supplied literature source: Groisser--Jung--Schwartzman, scaling--rotation curves, \url{https://arxiv.org/abs/1602.01187}.
\item Supplied literature source: Scaling--rotation uniqueness questions, \url{https://arxiv.org/abs/1702.03237}.
\item Supplied literature source: Oberwolfach Report 3/2018, \url{https://doi.org/10.4171/owr/2018/3}.
\item Literature assessment dated exactly \texttt{2026-08-21}: the supplied curation labels the problem open and states that low-dimensional smooth scaling--rotation curves are understood in several cases, while the requested algebraic characterization of shortest curves among all piecewise-smooth curves remains open. This is supplied dated curation derived from the OWR material, not an independently verified literature assessment; its status and cited claims require independent verification.
\item The supplied statement status is \texttt{corrected\_verified}; its accompanying explanation says that notation, typography, and sentence structure were repaired while preserving mathematical content. This is a supplied editorial label, not mathematical verification of this draft.
\end{itemize}
\end{document}
